% PAGES: 269-269
% Continues the sentence ending at the bottom of sec09_c1_2.tex (folio 252):
% "Let $\rho_{j,k}$ be the correlation between the" is immediately followed
% here by "returns $R_j$ and $R_k$..." (folio 253). No LaTeX construct was
% left open across the boundary.

returns $R_j$ and $R_k$, for $1 \le j \ne k \le n$, and let $\Sigma_R$ be the $n \times
n$ covariance matrix of the returns of the $n$ assets given by
\begin{align*}
\Sigma_R(j,k) &= \cov(R_j,R_k) \;=\; \sigma_j\sigma_k\rho_{j,k}, \quad \forall\, 1\le
j\ne k\le n; \\
\Sigma_R(i,i) &= \var(R_i) \;=\; \sigma_i^2, \quad \forall\, i = 1:n.
\end{align*}

Since the risk--free rate $r_f$ is a constant, we obtain from (9.2) that
\[
\sigma_R^2 \;=\; \var(R) \;=\; \var\!\left(\sum_{i=1}^n w_iR_i\right),
\]
and, using (7.25), we find that the variance of the return of the portfolio can be
written in matrix formulation as follows:
\begin{equation}
\sigma_R^2 \;=\; w^t\Sigma_Rw.
\end{equation}

A natural assumption we make throughout this chapter is that the return of none of the
assets can be replicated by using the other $n-1$ assets and cash.\footnote{Note the
similarity with the non--redundancy of the securities from a one period market model;
see Chapter 3.} In other words, we assume that there is no asset whose return is a
linear combination of the returns of the other $n-1$ assets and the risk--free return.
Note that, from Lemma 7.5, it follows that this assumption is equivalent to assuming
that the covariance matrix $\Sigma_R$ is nonsingular.

An important question is how should an efficient portfolio be set up?\footnote{Finding
efficient portfolios is one of the fundamental problems given a theoretical answer by
the modern portfolio theory of Markowitz and Sharpe; see Markowitz [28] and Sharpe [35]
for seminal papers.} In other words, what should the weights $w_i$, $i = 1:n$, of the
assets and the weight $w_{cash}$ of the cash position be in order to obtain a portfolio
(called a minimum variance portfolio) with the smallest variance of return, given a
fixed expected return, or to obtain a portfolio (called a maximum return portfolio)
with the highest expected return, given a fixed variance of the return?

\noindent\textbf{Minimum variance portfolio:}\\
\textit{Given $\mu_P$, find $w_{cash}$ and $w_i$, $i = 1:n$, with $\mu_R = \mu_P$, such
that $\sigma_R^2$ is minimal.}

\noindent\textbf{Maximum return portfolio:}\\
\textit{Given $\sigma_P$, find $w_{cash}$ and $w_i$, $i = 1:n$, with $\sigma_R^2 =
\sigma_P^2$, such that $\mu_R$ is maximal.}

In section 9.2, we include the solutions to the minimum variance portfolio problem and
to the maximum return portfolio problem, which will be derived in sections 9.3 and 9.4,
as well as the tangency portfolio formulation of these solutions. The pseudocodes for
finding the efficient portfolio allocation for both the minimum variance portfolio
problem and the maximum return portfolio are also included in section 9.2.

Note that the variances of asset returns and the covariance matrix for asset returns
can be estimated from historical data. However, estimated expected returns accurately
from historical is not feasible; see, e.g., Fabozzi and Markowitz [14] for alternative
approaches to finding asset allocations for efficient portfolios.

We conclude this section by proving formula (9.2).

Let $V(0)$ be the value at time 0 of a portfolio invested in $n$ assets and with a cash
position. Let $V_i(0)$ be the value of the portfolio position in asset $i$ at time 0, for
